\documentclass[12pt]{article}

% Idioma e codificação
\usepackage[brazil]{babel}
\usepackage[utf8]{inputenc}
\usepackage[T1]{fontenc}

% Matemática e listas
\usepackage{amsmath, amssymb}
\usepackage{enumitem}

% Layout e cores
\usepackage{geometry}
\usepackage{graphicx}
\usepackage{xcolor}

% Cores institucionais
\definecolor{maincolor}{RGB}{0,100,50}
\definecolor{forestgreen}{RGB}{0,100,50}

% Cabeçalho e rodapé
\usepackage{fancyhdr}

% Blocos
\usepackage[most]{tcolorbox}
\newtcolorbox{block}[1]{%
  colback=green!5,
  colframe=maincolor,
  coltitle=white,
  title=#1,
  fonttitle=\bfseries,
  boxrule=1.5pt,
  arc=3mm,
  left=6pt,
  right=6pt,
  top=6pt,
  bottom=6pt
}

% Margens
\geometry{margin=2.5cm}

% Fancyhdr
\pagestyle{fancy}
\fancyhf{}

\renewcommand{\headrulewidth}{2pt}
\renewcommand{\footrulewidth}{2pt}

\renewcommand{\headrule}{\hbox to\headwidth{%
  \color{forestgreen}\leaders\hrule height \headrulewidth\hfill}}

\renewcommand{\footrule}{\hbox to\headwidth{%
  \color{forestgreen}\leaders\hrule height \footrulewidth\hfill}}

% Altura do cabeçalho
\setlength{\headheight}{52pt}
\addtolength{\topmargin}{-20pt}

% Cabeçalho
\fancyhead[L]{%
  % Substitua pelo caminho correto da imagem ou mantenha o texto se não houver arquivo
  \textbf{UniRV} \\ \small Universidade de Rio Verde
}

\fancyhead[R]{%
  \raggedleft
  \textcolor{forestgreen}{%
    \small
    \textbf{Lista de Exercícios: Arquiteturas Tolerantes}\\
    Me.~Sergio Souza Novak\\
    \textit{Engenharia da Qualidade e Confiabilidade}
    }
}
    
% Rodapé
\fancyfoot[C]{\textcolor{forestgreen}{\thepage}}

% Comando para caixas de resposta
\newtcolorbox{resposta}[1]{%
  colback=white,
  colframe=forestgreen!30,
  arc=1mm,
  height=#1,
  width=\textwidth,
  boxrule=0.8pt,
  breakable
}

% =========================
\begin{document}
% =========================

  \begin{center}
    {\Large \textbf{Lista de Exercícios: Arquiteturas e Confiabilidade}}\\[0.2cm]
    \textit{Engenharia de Software / Engenharia da Qualidade}
  \end{center}

  \vspace{0.4cm}

  \section*{Parte 1 — Conceitos Fundamentais}

  \textbf{1.} Explique o objetivo principal de uma arquitetura tolerante a defeitos.
  \begin{resposta}{2.5cm}
  \end{resposta}

  \textbf{2.} Diferencie Prevenção, Detecção e Tolerância a defeitos.
  \begin{resposta}{3.5cm}
  \end{resposta}

  \textbf{3.} Por que sistemas críticos não podem depender apenas de testes para garantir segurança?
  \begin{resposta}{2.5cm}
  \end{resposta}

  \textbf{4.} Um sistema pode conter defeitos e ainda ser confiável? Justifique.
  \begin{resposta}{2.5cm}
  \end{resposta}

  \newpage

  \section*{Parte 2 — Hadoop e Sistemas Distribuídos}

  \textbf{5.} Explique por que o Hadoop foi projetado assumindo que falhas de hardware são normais.
  \begin{resposta}{2.5cm}
  \end{resposta}

  \textbf{6.} Descreva a função de cada componente: HDFS, MapReduce e YARN.
  \begin{resposta}{4.5cm}
  \end{resposta}

  \textbf{7.} Qual é a função do \textit{heartbeat} em um cluster?
  \begin{resposta}{8cm}
  \end{resposta}

  \textbf{8.} O que acontece quando um nó do cluster falha e como o fator de replicação (ex: 3) protege os dados?
  \begin{resposta}{8cm}
  \end{resposta}

  \section*{Parte 3 — Sistemas de Proteção}

  \textbf{9.} O que caracteriza um sistema de proteção e por que ele deve ser mais simples que o sistema principal?
  \begin{resposta}{3.5cm}
  \end{resposta}

  \textbf{10.} Explique o conceito matemático de probabilidade de falha sob demanda: $P(\text{falha} \mid \text{demanda})$.
  \begin{resposta}{2.5cm}
  \end{resposta}

  \newpage

  \section*{Parte 4 — Métricas Quantitativas}

  \begin{block}{Fórmulas de Referência}
      \small
      \begin{itemize}
          \item $POFOD = \frac{\text{Nº de Falhas}}{\text{Nº de Solicitações}}$ \quad | \quad $ROCOF = \frac{\text{Nº de Falhas}}{\text{Tempo Total}}$
          \item $MTTF = \frac{1}{ROCOF}$ \quad | \quad $MTTR = \frac{\text{Tempo Total de Reparo}}{\text{Nº de Falhas}}$
          \item $AVAIL = \frac{MTTF}{MTTF + MTTR}$
      \end{itemize}
  \end{block}

  \textbf{11.} Cada canal de um sistema possui $POFOD = 0,001$. Calcule a confiabilidade do sistema com três canais independentes (redundância).
  \begin{resposta}{3.5cm}
  \end{resposta}

  \textbf{12.} Para atingir $99,99\%$ de disponibilidade com $MTTF = 4.000h$, qual deve ser o $MTTR$ máximo (em minutos)?
  \begin{resposta}{3.5cm}
  \end{resposta}

  \section*{Parte 5 — Diversidade e N-Versões}

  \textbf{13.} O que é Programação N-versões e por que $N \ge 3$ para tolerar uma falha?
  \begin{resposta}{8cm}
  \end{resposta}

  \textbf{14.} Explique o conceito de "Falha de Modo Comum" e como a diversidade de hardware/software tenta mitigá-la.
  \begin{resposta}{3.5cm}
  \end{resposta}

  \textbf{15.} Quando Confiabilidade é mais importante que Disponibilidade? Dê um exemplo prático.
  \begin{resposta}{3.5cm}
  \end{resposta}

  \vspace{1cm}

  \begin{center}
      \textit{“Boa atividade! Att. Me. Prof. Sergio Souza Novak.”}
  \end{center}

% =========================
\end{document}
% =========================